\subsection{在拓扑空间上自由作用的群}

定义 2. 设 G 是作用于一个集合 X 上的群。如果 X 的每个元素的稳定子群都是 \(\{e\}\)，换句话说，如果关系 \(s.x = x, s \in G, x \in X\) 蕴含 s = e，则称 G 自由作用于 X。

例. 设 G 是群，H 是 G 的子群。则 H 通过（左或右）平移自由作用于 G。

设 G 是自由作用于集合 X 的群，R 是 G 在 X 上定义的等价关系，\(C \subset X \times X\) 是 R 的图。若 \((x, y) \in C\)，则存在 \(s \in G\) 使得 s.x = y；且 s 是唯一的，因为 \(s.x = s'.x\) 蕴含 \(s'^{-1}s.x = x\)，从而 \(s'^{-1}s = e\)（由于 G 自由作用）。让 \((x, y) \in C\) 对应于使 s.x = y 的唯一的 \(s \in G\)，我们就定义了一个映射 \(\varphi: C \to G\)，我们称它为 C 到 G 中的典范映射。采用此记号：

命题 6. 设 G 是连续作用于拓扑空间 X 的拓扑群，且设 G 自由作用于 X。则 G 逆紧作用于 X 当且仅当下述条件成立：

(FP) 由 G 所定义的等价关系的图 C 在 \(X \times X\) 中闭，且典范映射 \(\varphi: C \rightarrow G\) 连续。

集合 C 是映射 \(\theta: (s, x) \to (x, s.x)\)（\(\mathbf{G} \times \mathbf{X}\) 到 \(\mathbf{X} \times \mathbf{X}\) 中的映射）的像。我们知道（第一章，§ 10，第 1 节，命题 2），\(\theta\) 是逆紧的当且仅当 C 在 \(\mathbf{X} \times \mathbf{X}\) 中闭，并且（若用 \(\theta'\) 表示把 \(\theta\) 看作 \(\mathbf{G} \times \mathbf{X}\) 到 C 中的映射）\(\theta'\) 是同胚。现在，假设蕴含 \(\theta'\) 是双射且其逆是映射 \((x, y) \to (\varphi(x, y), x)\)。于是 \(\theta'\) 是同胚当且仅当 \(\varphi\) 连续。

